\chapter{Introduction}

Citizen-science bird monitoring generates observations at scales impossible for
professional ornithologists alone.
Platforms such as eBird receive hundreds of millions of observations annually,
yet data quality depends critically on observers' ability to identify species in
the field~\cite{Sullivan2009}.
Mobile applications with automated species recognition lower this barrier, but
most existing tools either require continuous network access, cover limited
geographic ranges, or rely on a single sensing modality.

BirdNET~\cite{BirdNET2021} demonstrated that a convolutional network trained on
spectrograms can achieve macro-averaged AUC $> 0.99$ across thousands of species,
making robust audio-based identification practical on consumer hardware.
Image-based classifiers such as MobileNetV2~\cite{Sandler2018} offer a
complementary modality; the two are rarely combined in an offline mobile context.

LogChirpy addresses three concrete gaps:

\begin{enumerate}
  \item \textbf{Multi-modal offline identification.}
        Combining BirdNET audio and a MobileNetV2 image classifier on-device,
        without any cloud inference dependency.

  \item \textbf{Large-scale multilingual reference database.}
        The Clements 2024 checklist~\cite{Clements2024} with 33,241 species
        in six languages and 9,331 Wikimedia Commons reference images, fully
        accessible offline.

  \item \textbf{Sequential pipeline architecture.}
        A unified service layer that eliminates the file-lock race conditions
        observed when audio recording and camera capture run concurrently on
        Android.
\end{enumerate}

The remainder of this paper describes the system architecture (Chapter~2),
an independent benchmark of the image classifier (Chapter~3), and a
discussion of limitations and scope (Chapter~4).
